% Curs 12 - Pagina 6 (3-I-1995).
% Fragment pentru % Curs 12 - Pagina 6 (3-I-1995).
% Fragment pentru % Curs 12 - Pagina 6 (3-I-1995).
% Fragment pentru % Curs 12 - Pagina 6 (3-I-1995).
% Fragment pentru \input{courses/pages/curs12-pagina6-ro}.
% Fara preambul sau salt de pagina; ordinea manuscrisului este pastrata.
% NOTE EDITORIALE:
% - La inceput, afirmatia R* subset C necesita ca C sa fie o congruenta;
%   in manuscris C este numita doar submultime a lui A* x A*.
% - Definitia lui R de aici (x,y litere din A, x != y) este mai restransa
%   decat relatia folosita in Aplicatia 1, la pagina 5.
% - Prima egalitate care mentioneaza generatorii A* si A-bar este putin
%   neclara; aici A* = <A> si A*/R* = <A-bar> exprima lectura contextuala.
% - In demonstratia punctului (2) apar sageti/semne de intrebare; sunt
%   pastrate, fara a completa finalul demonstratiei din pagina urmatoare.

\subsection*{Curs 12 - Pagina 6}

\noindent $R^*$ e tranzitiv\u{a}.
\[
R\subseteq S_R\subseteq R^*.
\]

\noindent Fie $C\subseteq A^*\times A^*$ a.\^{\i}.\ $R\subseteq C$
$\Rightarrow R^*\subseteq C$.

\bigskip
\noindent\underline{Proprietatea de universalitate a monoidului comutativ generat de $A$}

\medskip
\noindent\fbox{T}\quad Fie $A\neq\varnothing$,
\[
R=\bigl\{(xy,yx)\in A^*\times A^*\mid x,y\in A,\ x\neq y\bigr\},
\]
\noindent \c{s}i $R^*$ congruen\c{t}a pe $A^*$ generat\u{a} de $R$.

\medskip
\noindent\textbf{1.} Monoidul $\dfrac{A^*}{R^*}$ e comutativ.

\medskip
\noindent\textbf{2.} Aplica\c{t}ia $\eta_A:A\longrightarrow\dfrac{A^*}{R^*}$,
\[
\eta_A=\Bigl(A^*\xrightarrow{\ \rho\ }\frac{A^*}{R^*}\Bigr)
\circ\bigl(A\xrightarrow{\ \zeta_A\ }A^*\bigr),
\]
\noindent $A\xrightarrow{\eta_A}\dfrac{A^*}{R^*}$, pentru orice
$(M,\cdot,e)$ monoid comutativ \c{s}i orice $\eta:A\to M$,
\[
\exists!\,g:\frac{A^*}{R^*}\longrightarrow M
\quad\text{a.\^{\i}.}\quad g\circ\eta_A=\eta.
\]

\bigskip
\noindent\underline{Dem.}

\medskip
\noindent\textbf{(1)} $A\subseteq A^*$.
\[
A^*=\langle A\rangle\quad\Longrightarrow\quad
\frac{A^*}{R^*}=\langle\overline A\rangle,
\qquad \overline A=\{\overline x\mid x\in A\},
\quad\overline x\in\frac{A^*}{R^*}.
\]

\noindent Ar\u{a}t\u{a}m c\u{a}
\[
\overline x\,\overline y=\overline y\,\overline x
\qquad (\forall)\,x,y\in A,
\]
\noindent deoarece $(xy,yx)\in R$.
\[
\overline{xy}=\overline{yx},\qquad
\overline u\,\overline v=\overline{uv}
\quad (\forall)\,\overline u,\overline v\in\frac{A^*}{R^*}.
\]

\medskip
\noindent\textbf{(2)}
\[
\begin{tikzcd}[column sep=large, row sep=large]
A \arrow[r,"\zeta_A"] \arrow[dr,"\eta"']
  \arrow[rr,bend left=28,"\eta_A"]
  & A^* \arrow[r,"\rho"] \arrow[d,dashed,"{\exists!\,f}"']
  & \dfrac{A^*}{R^*} \arrow[dl,dashed,"{\exists!\,g}"] \\
& M &
\end{tikzcd}
\]

\noindent $\exists!\,f$ a.\^{\i}.\ $f\circ\zeta_A=\eta$.
\[
R^*\subseteq R_f\quad\Longleftarrow\quad R\subseteq R_f.
\]
\noindent $(xy,yx)\in R$,
\[
f(xy)\overset{?}{=}f(yx)
\quad\Longrightarrow\quad (xy,yx)\in R_f.
\]
\[
f(xy)=f(x)f(y)=f(y)f(x),
\qquad f\circ\zeta_A=\eta,\quad g\circ\rho=f,\quad
 g\circ\eta_A=\eta.
\]
% Ultima linie a demonstratiei este foarte compacta in manuscris;
% egalitatile de compunere de la final sunt citite contextual.
.
% Fara preambul sau salt de pagina; ordinea manuscrisului este pastrata.
% NOTE EDITORIALE:
% - La inceput, afirmatia R* subset C necesita ca C sa fie o congruenta;
%   in manuscris C este numita doar submultime a lui A* x A*.
% - Definitia lui R de aici (x,y litere din A, x != y) este mai restransa
%   decat relatia folosita in Aplicatia 1, la pagina 5.
% - Prima egalitate care mentioneaza generatorii A* si A-bar este putin
%   neclara; aici A* = <A> si A*/R* = <A-bar> exprima lectura contextuala.
% - In demonstratia punctului (2) apar sageti/semne de intrebare; sunt
%   pastrate, fara a completa finalul demonstratiei din pagina urmatoare.

\subsection*{Curs 12 - Pagina 6}

\noindent $R^*$ e tranzitiv\u{a}.
\[
R\subseteq S_R\subseteq R^*.
\]

\noindent Fie $C\subseteq A^*\times A^*$ a.\^{\i}.\ $R\subseteq C$
$\Rightarrow R^*\subseteq C$.

\bigskip
\noindent\underline{Proprietatea de universalitate a monoidului comutativ generat de $A$}

\medskip
\noindent\fbox{T}\quad Fie $A\neq\varnothing$,
\[
R=\bigl\{(xy,yx)\in A^*\times A^*\mid x,y\in A,\ x\neq y\bigr\},
\]
\noindent \c{s}i $R^*$ congruen\c{t}a pe $A^*$ generat\u{a} de $R$.

\medskip
\noindent\textbf{1.} Monoidul $\dfrac{A^*}{R^*}$ e comutativ.

\medskip
\noindent\textbf{2.} Aplica\c{t}ia $\eta_A:A\longrightarrow\dfrac{A^*}{R^*}$,
\[
\eta_A=\Bigl(A^*\xrightarrow{\ \rho\ }\frac{A^*}{R^*}\Bigr)
\circ\bigl(A\xrightarrow{\ \zeta_A\ }A^*\bigr),
\]
\noindent $A\xrightarrow{\eta_A}\dfrac{A^*}{R^*}$, pentru orice
$(M,\cdot,e)$ monoid comutativ \c{s}i orice $\eta:A\to M$,
\[
\exists!\,g:\frac{A^*}{R^*}\longrightarrow M
\quad\text{a.\^{\i}.}\quad g\circ\eta_A=\eta.
\]

\bigskip
\noindent\underline{Dem.}

\medskip
\noindent\textbf{(1)} $A\subseteq A^*$.
\[
A^*=\langle A\rangle\quad\Longrightarrow\quad
\frac{A^*}{R^*}=\langle\overline A\rangle,
\qquad \overline A=\{\overline x\mid x\in A\},
\quad\overline x\in\frac{A^*}{R^*}.
\]

\noindent Ar\u{a}t\u{a}m c\u{a}
\[
\overline x\,\overline y=\overline y\,\overline x
\qquad (\forall)\,x,y\in A,
\]
\noindent deoarece $(xy,yx)\in R$.
\[
\overline{xy}=\overline{yx},\qquad
\overline u\,\overline v=\overline{uv}
\quad (\forall)\,\overline u,\overline v\in\frac{A^*}{R^*}.
\]

\medskip
\noindent\textbf{(2)}
\[
\begin{tikzcd}[column sep=large, row sep=large]
A \arrow[r,"\zeta_A"] \arrow[dr,"\eta"']
  \arrow[rr,bend left=28,"\eta_A"]
  & A^* \arrow[r,"\rho"] \arrow[d,dashed,"{\exists!\,f}"']
  & \dfrac{A^*}{R^*} \arrow[dl,dashed,"{\exists!\,g}"] \\
& M &
\end{tikzcd}
\]

\noindent $\exists!\,f$ a.\^{\i}.\ $f\circ\zeta_A=\eta$.
\[
R^*\subseteq R_f\quad\Longleftarrow\quad R\subseteq R_f.
\]
\noindent $(xy,yx)\in R$,
\[
f(xy)\overset{?}{=}f(yx)
\quad\Longrightarrow\quad (xy,yx)\in R_f.
\]
\[
f(xy)=f(x)f(y)=f(y)f(x),
\qquad f\circ\zeta_A=\eta,\quad g\circ\rho=f,\quad
 g\circ\eta_A=\eta.
\]
% Ultima linie a demonstratiei este foarte compacta in manuscris;
% egalitatile de compunere de la final sunt citite contextual.
.
% Fara preambul sau salt de pagina; ordinea manuscrisului este pastrata.
% NOTE EDITORIALE:
% - La inceput, afirmatia R* subset C necesita ca C sa fie o congruenta;
%   in manuscris C este numita doar submultime a lui A* x A*.
% - Definitia lui R de aici (x,y litere din A, x != y) este mai restransa
%   decat relatia folosita in Aplicatia 1, la pagina 5.
% - Prima egalitate care mentioneaza generatorii A* si A-bar este putin
%   neclara; aici A* = <A> si A*/R* = <A-bar> exprima lectura contextuala.
% - In demonstratia punctului (2) apar sageti/semne de intrebare; sunt
%   pastrate, fara a completa finalul demonstratiei din pagina urmatoare.

\subsection*{Curs 12 - Pagina 6}

\noindent $R^*$ e tranzitiv\u{a}.
\[
R\subseteq S_R\subseteq R^*.
\]

\noindent Fie $C\subseteq A^*\times A^*$ a.\^{\i}.\ $R\subseteq C$
$\Rightarrow R^*\subseteq C$.

\bigskip
\noindent\underline{Proprietatea de universalitate a monoidului comutativ generat de $A$}

\medskip
\noindent\fbox{T}\quad Fie $A\neq\varnothing$,
\[
R=\bigl\{(xy,yx)\in A^*\times A^*\mid x,y\in A,\ x\neq y\bigr\},
\]
\noindent \c{s}i $R^*$ congruen\c{t}a pe $A^*$ generat\u{a} de $R$.

\medskip
\noindent\textbf{1.} Monoidul $\dfrac{A^*}{R^*}$ e comutativ.

\medskip
\noindent\textbf{2.} Aplica\c{t}ia $\eta_A:A\longrightarrow\dfrac{A^*}{R^*}$,
\[
\eta_A=\Bigl(A^*\xrightarrow{\ \rho\ }\frac{A^*}{R^*}\Bigr)
\circ\bigl(A\xrightarrow{\ \zeta_A\ }A^*\bigr),
\]
\noindent $A\xrightarrow{\eta_A}\dfrac{A^*}{R^*}$, pentru orice
$(M,\cdot,e)$ monoid comutativ \c{s}i orice $\eta:A\to M$,
\[
\exists!\,g:\frac{A^*}{R^*}\longrightarrow M
\quad\text{a.\^{\i}.}\quad g\circ\eta_A=\eta.
\]

\bigskip
\noindent\underline{Dem.}

\medskip
\noindent\textbf{(1)} $A\subseteq A^*$.
\[
A^*=\langle A\rangle\quad\Longrightarrow\quad
\frac{A^*}{R^*}=\langle\overline A\rangle,
\qquad \overline A=\{\overline x\mid x\in A\},
\quad\overline x\in\frac{A^*}{R^*}.
\]

\noindent Ar\u{a}t\u{a}m c\u{a}
\[
\overline x\,\overline y=\overline y\,\overline x
\qquad (\forall)\,x,y\in A,
\]
\noindent deoarece $(xy,yx)\in R$.
\[
\overline{xy}=\overline{yx},\qquad
\overline u\,\overline v=\overline{uv}
\quad (\forall)\,\overline u,\overline v\in\frac{A^*}{R^*}.
\]

\medskip
\noindent\textbf{(2)}
\[
\begin{tikzcd}[column sep=large, row sep=large]
A \arrow[r,"\zeta_A"] \arrow[dr,"\eta"']
  \arrow[rr,bend left=28,"\eta_A"]
  & A^* \arrow[r,"\rho"] \arrow[d,dashed,"{\exists!\,f}"']
  & \dfrac{A^*}{R^*} \arrow[dl,dashed,"{\exists!\,g}"] \\
& M &
\end{tikzcd}
\]

\noindent $\exists!\,f$ a.\^{\i}.\ $f\circ\zeta_A=\eta$.
\[
R^*\subseteq R_f\quad\Longleftarrow\quad R\subseteq R_f.
\]
\noindent $(xy,yx)\in R$,
\[
f(xy)\overset{?}{=}f(yx)
\quad\Longrightarrow\quad (xy,yx)\in R_f.
\]
\[
f(xy)=f(x)f(y)=f(y)f(x),
\qquad f\circ\zeta_A=\eta,\quad g\circ\rho=f,\quad
 g\circ\eta_A=\eta.
\]
% Ultima linie a demonstratiei este foarte compacta in manuscris;
% egalitatile de compunere de la final sunt citite contextual.
.
% Fara preambul sau salt de pagina; ordinea manuscrisului este pastrata.
% NOTE EDITORIALE:
% - La inceput, afirmatia R* subset C necesita ca C sa fie o congruenta;
%   in manuscris C este numita doar submultime a lui A* x A*.
% - Definitia lui R de aici (x,y litere din A, x != y) este mai restransa
%   decat relatia folosita in Aplicatia 1, la pagina 5.
% - Prima egalitate care mentioneaza generatorii A* si A-bar este putin
%   neclara; aici A* = <A> si A*/R* = <A-bar> exprima lectura contextuala.
% - In demonstratia punctului (2) apar sageti/semne de intrebare; sunt
%   pastrate, fara a completa finalul demonstratiei din pagina urmatoare.

\subsection*{Curs 12 - Pagina 6}

\noindent $R^*$ e tranzitiv\u{a}.
\[
R\subseteq S_R\subseteq R^*.
\]

\noindent Fie $C\subseteq A^*\times A^*$ a.\^{\i}.\ $R\subseteq C$
$\Rightarrow R^*\subseteq C$.

\bigskip
\noindent\underline{Proprietatea de universalitate a monoidului comutativ generat de $A$}

\medskip
\noindent\fbox{T}\quad Fie $A\neq\varnothing$,
\[
R=\bigl\{(xy,yx)\in A^*\times A^*\mid x,y\in A,\ x\neq y\bigr\},
\]
\noindent \c{s}i $R^*$ congruen\c{t}a pe $A^*$ generat\u{a} de $R$.

\medskip
\noindent\textbf{1.} Monoidul $\dfrac{A^*}{R^*}$ e comutativ.

\medskip
\noindent\textbf{2.} Aplica\c{t}ia $\eta_A:A\longrightarrow\dfrac{A^*}{R^*}$,
\[
\eta_A=\Bigl(A^*\xrightarrow{\ \rho\ }\frac{A^*}{R^*}\Bigr)
\circ\bigl(A\xrightarrow{\ \zeta_A\ }A^*\bigr),
\]
\noindent $A\xrightarrow{\eta_A}\dfrac{A^*}{R^*}$, pentru orice
$(M,\cdot,e)$ monoid comutativ \c{s}i orice $\eta:A\to M$,
\[
\exists!\,g:\frac{A^*}{R^*}\longrightarrow M
\quad\text{a.\^{\i}.}\quad g\circ\eta_A=\eta.
\]

\bigskip
\noindent\underline{Dem.}

\medskip
\noindent\textbf{(1)} $A\subseteq A^*$.
\[
A^*=\langle A\rangle\quad\Longrightarrow\quad
\frac{A^*}{R^*}=\langle\overline A\rangle,
\qquad \overline A=\{\overline x\mid x\in A\},
\quad\overline x\in\frac{A^*}{R^*}.
\]

\noindent Ar\u{a}t\u{a}m c\u{a}
\[
\overline x\,\overline y=\overline y\,\overline x
\qquad (\forall)\,x,y\in A,
\]
\noindent deoarece $(xy,yx)\in R$.
\[
\overline{xy}=\overline{yx},\qquad
\overline u\,\overline v=\overline{uv}
\quad (\forall)\,\overline u,\overline v\in\frac{A^*}{R^*}.
\]

\medskip
\noindent\textbf{(2)}
\[
\begin{tikzcd}[column sep=large, row sep=large]
A \arrow[r,"\zeta_A"] \arrow[dr,"\eta"']
  \arrow[rr,bend left=28,"\eta_A"]
  & A^* \arrow[r,"\rho"] \arrow[d,dashed,"{\exists!\,f}"']
  & \dfrac{A^*}{R^*} \arrow[dl,dashed,"{\exists!\,g}"] \\
& M &
\end{tikzcd}
\]

\noindent $\exists!\,f$ a.\^{\i}.\ $f\circ\zeta_A=\eta$.
\[
R^*\subseteq R_f\quad\Longleftarrow\quad R\subseteq R_f.
\]
\noindent $(xy,yx)\in R$,
\[
f(xy)\overset{?}{=}f(yx)
\quad\Longrightarrow\quad (xy,yx)\in R_f.
\]
\[
f(xy)=f(x)f(y)=f(y)f(x),
\qquad f\circ\zeta_A=\eta,\quad g\circ\rho=f,\quad
 g\circ\eta_A=\eta.
\]
% Ultima linie a demonstratiei este foarte compacta in manuscris;
% egalitatile de compunere de la final sunt citite contextual.
